% ----------------------
\subsection{Section 28 (September 1830)}
% ----------------------
	\label{DC28}
	
	% -----
	\subsubsection{Historical Background}
	% -----
		\label{DC28-HB}
		
		% --
		\paragraph{JSP}
		% --
		
			``This revelation was a response to actions by Oliver Cowdery and Hiram Page that raised the question of whether JS was the only one authorized to deliver revelation to the church. The question first arose in summer 1830 when Oliver Cowdery `commanded' JS to change a passage in `\hyperref[DC20]{Articles and Covenants},' a document outlining the basic beliefs and practices of the Church of Christ. When the church was organized on 6 April 1830, Cowdery was ordained second elder while JS was ordained first elder, and the two worked closely together to oversee the newly formed organization. Presumably Cowdery worked with JS in preparing Articles and Covenants and was part of the `unanimous voice of the whole congregation' that approved Articles and Covenants at the first conference of the church on 9 June, yet weeks later he sent a letter ordering JS to make a correction to the document. According to JS's history, Cowdery objected to the requirement that candidates for baptism `truly manifest by their works that they have received the gift of Christ unto the remission of their sins,' and he wrote to JS, `I command you in the name of God to erase those words, that no priestcraft be amongst us.'%
				\footnote{%
					% \label{}%
					The source here is HC, and TCHC doesn't point to any other source for this remark of Cowdery.
					}
			In response, JS traveled from Harmony, Pennsylvania, to Fayette, New York, to persuade Cowdery and the Whitmers that they were mistaken. According to JS's later account, it was `not without both labor and perseverance' that he `could prevail with any of them to reason calmly on the subject.' Finally, with support from Christian Whitmer, JS convinced Cowdery and the Whitmer family `that they had been in error, and that the sentence in dispute was in accordance of the rest of the commandment.'
			
			``The second challenge to JS's authority came in early September, when JS and Emma Smith moved from Harmony to Fayette and found to their `great grief' that `Brother Hyrum [Hiram] Page had got in his possession, a certain stone, by which he had obtained to certain revelations...all of which were entirely at variance with the order of Gods house, as laid down in the new Testament, as well as in our late revelations.'%
				\footnote{%
					% \label{}%
					JSP note: ``JS History, vol. A-1, 53--54. Newel Knight recalled that Page `had quite a roll of papers full of these revelations, and many in the church were led astray by them.' Ezra Booth, who wrote a series of antagonistic letters denouncing JS after leaving the church in the fall of 1831, explained his understanding of Page's seer stone: `[He] found a smooth stone, upon which there appeared to be writing, which when transcribed upon paper, disappeared from the stone, and another impression appeared in its place. This when copied, vanished as the former had done, and so it continued alternately appearing and disappearing; in the meanwhile, he continued to write, until he had written over considerable paper. It bore most striking marks of a Mormonite revelation, and was received as an authentic document by most of the Mormonites, till Smith, by his superior sagacity, discovered it to be a Satanic fraud.' George A. Smith later stated that the stone was black and explained that on it Page saw `certain characters' that he copied down as revelations. Emer Harris also recalled that Page's stone was black; he added that it was destroyed. (Knight, History, 146; `Letters from David and John C. Whitmer,' \textit{Saints' Herald}, 5 Feb. 1887, 90; Ezra Booth, `Mormonism---Nos. VIII--IX,' \textit{Ohio Star} [Ravenna], 8 Dec. 1831, [1]; George A. Smith, in \textit{Journal of Discourses}, 15 Nov. 1864, 11:2; Provo, UT, Central Stake, General Minutes, 6 Apr. 1856, vol. 10, p. 273.)''
					}
			With another conference of the church approaching, JS initially `thought it wisdom not to do much more than to converse with the brethren on the subject, untill the conference should meet.' But upon finding that many, including Cowdery and the Whitmer family, supported Page, JS (apparently with Cowdery's encouragement) decided it would be `best to enquire of the Lord concerning so important a matter.' Before the conference convened, JS dictated the revelation featured here, which addressed the issues surrounding both Cowdery's role and `the things set forth by this [Page's] stone.'
			
			``While affirming that Cowdery was called to teach and would receive revelation, the text stated that he was not to write revelation to the church by commandment and added, `Thou shalt not command him which is at thy head \& at the head of the Church for I have given him the keys of the mysteries of the Revelations which are sealed until I shall appoint unto him another in his stead.' The revelation also directed Cowdery to explain to Hiram Page privately that the latter's revelations were from Satan. The role for JS described herein was ratified at the subsequent conference, convened 26 September, at which JS was appointed `by the voice of the Conference to receive and write Revelations \& Commandments for this Church.' Cowdery then read aloud Articles and Covenants, which he had previously criticized, and JS delivered comments on that document.
			
			``This revelation also called Cowdery to preach to the `Lamanites,' or American Indians, giving specificity to a July revelation for Cowdery that commanded him in general terms to preach the gospel. In addition, the text featured here commanded JS to establish a church among the Lamanites, where `the City [the New Jerusalem] shall be built.' Although the revelation declared that `no man' yet knew the location for the New Jerusalem, Cowdery signed a statement on 17 October declaring that he would preach and `rear up a pillar as a witness where the Temple of God shall be built, in the glorious New-Jerusalem.' A revelation earlier that month had affirmed that church elders would establish the New Jerusalem in America as an apocalyptic fulfillment of both biblical and Book of Mormon prophecy. After other men were called to accompany Cowdery, the group of missionaries left in late October 1830.''
			
		% --
		\paragraph{HDDC}
		% --
		
			``Since Newel Knight was an eye-witness to these events, the entry in his diary is valuable:
			
				\begin{quote}
					After arranging my affairs at home, I again set out for Fayette, to attend our second conference, which had been appointed to be held at Father Whitmer's where Joseph then resided. On my arrival I found Brother Joseph in great distress of mind on account of Hyrum Page, who had managed to get up some dissention of feeling among the brethren by giving revelations concerning the government of the Church and other matters, which he claimed to have received through the medium of a stone he possessed. He had quite a roll of papers full of these revelations, and many in the Church were led astray by them. Even Oliver Cowdery and the Whitmer family had given heed to them, although they were in contradiction to the New Testament and the revelations of these last days. Here was a chance for Satan to work among the little flock, and he sought by this means to accomplish what persecution failed to do. Joseph was perplexed and scarcely knew how to meet this new exigency. That night I occupied the same room that he did and the greater part of the night was spent in prayer and supplication. After much labor with these brethren they were convinced of their error, and confessed the same, renouncing the revelations as not being of God, but acknowledged that Satan had conspired ot overthrow their belief in the true plan of salvation. In consequence of these things Joseph enquired of the Lord before conference commenced and received the revelation published on page 140 of the Doctrine and Covenants, wherein God explicitly states His mind and will concerning the receiving of revelations.
					
					Conference having assembled, the first thing done was to consider the subject of the stone in connection with Hyrum Page, and after considerable investigation and discussion, Brother Page and all the members of the Church present renounced the stone, and the revelations connected with it, much to our joy and satisfaction.%
						\footnote{%
							% \label{}%
							HDDC note: ``Journal History of the Church, September 26, 1830, located in the HDC.'' It's PDF 98 of the 1830 JH, and like the quotation from JH in the HDDC background to \hyperref[DC27]{DC 27}, it's a cut-out from some other source. See my notes on section 27.
							}
				\end{quote}
				
				``...Even though Hiram Page's stone was renounced by its owner and the membership of the Church, it was kept as a souvenir. It eventually was placed in the hands of the leaders of the Reorganized Church of Jesus Christ of Latter Day Saints, and is now housed in their Department of History. Elder Cecil McGavin has had the privilege of examing [sic] this stone, and wrote the following description of it:
				
					\begin{quote}
						It is a flat stone about seven inches long, four wide, and one-quarter inch in thickness. It is dark gray in color with waves of brown and purple gracefully interwoven across the surface. A small hole has been drilled through one end of it as if a string had been threaded through it.%
							\footnote{%
								% \label{}%
								HDDC note: ``E. Cecil McGavin, \uline{The Historical Background of The Doctrine and Covenants}, First edition (Salt Lake City: Paragon Printing Co.), p. 93.''
								}
					\end{quote}
					
				``An interesting sidelight concerning this stone and the revelations received through it is found in a statement made by David Whitmer fifty-seven years later. He wrote:
				
					\begin{quote}
						As to the revelations which came through Hiram Page's stone, I will state that Oliver and I never thought much about them. We talked of them, and thought they might be from God, or might be from Satan.''%
						\footnote{%
							% \label{}%
							HDDC note: ``The Saints' Herald [Lamoni, Iowa], February 5, 1887, p. 90.''
							} 
					\end{quote}
			
	% -----
	\subsubsection{Versions}
	% -----
		\label{DC28-V}
		
		See the \href{https://drive.google.com/file/d/1goAvNz8jS4R8slYfMG_db55Ty2z_vmgK/view?usp=sharing}{sections comparison} document.
	
		\begin{compactdesc}
			\item[JSP] \hyperref[EMS-1830-JS-SEP-1830-B]{September 1830}
			\item[BC] Chapter 30, 67--68
			\item[DC 1835] Section 51, 181--182
			\item[TS] 4:8 (1 March 1843), 119
			\item[MS] 4:10 (February 1844), 152--153
			\item[DC 1844] Section 50, 272--274
		\end{compactdesc}

	% -----
	\subsubsection{Section 28:1} % *DC28-1 
	% -----
		\label{DC28-1}

		BEHOLD, I say unto thee, Oliver, that it shall be given unto thee that thou shalt be heard by the church in all things whatsoever thou shalt teach them by the \hyperref[COMFORTER]{Comforter}, concerning the revelations and commandments which I have given.

		% \begin{framed}
		% \scriptsize
			% \begin{compactdesc}
				% \item[JSP] 
				% % \item[BC] 
				% % \item[]
			% \end{compactdesc}
		% \end{framed}

	% -----
	\subsubsection{Section 28:2} % *DC28-2 
	% -----
		\label{DC28-2}

		But, behold, verily, verily, I say unto thee, no one shall be appointed to receive commandments and revelations in this church excepting my servant Joseph Smith, Jun., for he receiveth them even as Moses.

		% \begin{framed}
		% \scriptsize
			% \begin{compactdesc}
				% \item[JSP] 
				% % \item[BC] 
				% % \item[]
			% \end{compactdesc}
		% \end{framed}

	% -----
	\subsubsection{Section 28:3} % *DC28-3 
	% -----
		\label{DC28-3}

		And thou shalt be obedient unto the things which I shall give unto him, even as Aaron, to declare faithfully the commandments and the revelations, with power and authority unto the church.

		% \begin{framed}
		% \scriptsize
			% \begin{compactdesc}
				% \item[JSP] 
				% % \item[BC] 
				% % \item[]
			% \end{compactdesc}
		% \end{framed}

	% -----
	\subsubsection{Section 28:4} % *DC28-4 
	% -----
		\label{DC28-4}

		And if thou art led at any time by the \hyperref[COMFORTER]{Comforter} to speak or teach, or at all times by the way of commandment unto the church, thou mayest do it.

		% \begin{framed}
		% \scriptsize
			% \begin{compactdesc}
				% \item[JSP] 
				% % \item[BC] 
				% % \item[]
			% \end{compactdesc}
		% \end{framed}

	% -----
	\subsubsection{Section 28:5} % *DC28-5 
	% -----
		\label{DC28-5}

		But thou shalt not write by way of commandment, but by wisdom;%
			\footnote{%
				% \label{}%
				This way of putting the instruction to Cowdery makes me think of \hyperref[DC89]{DC 89}, but also \hyperref[DC88-118]{DC 88:118}.
				}

		% \begin{framed}
		% \scriptsize
			% \begin{compactdesc}
				% \item[JSP] 
				% % \item[BC] 
				% % \item[]
			% \end{compactdesc}
		% \end{framed}

	% -----
	\subsubsection{Section 28:6} % *DC28-6 
	% -----
		\label{DC28-6}

		And thou shalt not command him who is at thy head, and at the head of the church;

		% \begin{framed}
		% \scriptsize
			% \begin{compactdesc}
				% \item[JSP] 
				% % \item[BC] 
				% % \item[]
			% \end{compactdesc}
		% \end{framed}

	% -----
	\subsubsection{Section 28:7} % *DC28-7 
	% -----
		\label{DC28-7}

		For I have given him the keys of the \hyperref[MYSTERY]{mysteries}, and the revelations which are sealed, until I shall appoint unto them another in his stead.

		% \begin{framed}
		% \scriptsize
			% \begin{compactdesc}
				% \item[JSP] 
				% % \item[BC] 
				% % \item[]
			% \end{compactdesc}
		% \end{framed}

	% -----
	\subsubsection{Section 28:8} % *DC28-8 
	% -----
		\label{DC28-8}

		And now, behold, I say unto you that you shall go unto the Lamanites and preach my gospel unto them; and inasmuch as they receive thy teachings thou shalt cause my church to be established among them; and thou shalt have revelations, but write them not by way of commandment.

		% \begin{framed}
		% \scriptsize
			% \begin{compactdesc}
				% \item[JSP] 
				% % \item[BC] 
				% % \item[]
			% \end{compactdesc}
		% \end{framed}

	% -----
	\subsubsection{Section 28:9} % *DC28-9 
	% -----
		\label{DC28-9}

		And now, behold, I say unto you that it is not revealed, and no man knoweth where the city Zion shall be built, but it shall be given hereafter.  Behold, I say unto you that it shall be on the borders by the Lamanites.

		% \begin{framed}
		% \scriptsize
			% \begin{compactdesc}
				% \item[JSP] 
				% % \item[BC] 
				% % \item[]
			% \end{compactdesc}
		% \end{framed}

	% -----
	\subsubsection{Section 28:10} % *DC28-10 
	% -----
		\label{DC28-10}

		Thou shalt not leave this place until after the conference; and my servant Joseph shall be appointed to preside over the conference by the voice of it, and what he saith to thee thou shalt tell.

		% \begin{framed}
		% \scriptsize
			% \begin{compactdesc}
				% \item[JSP] 
				% % \item[BC] 
				% % \item[]
			% \end{compactdesc}
		% \end{framed}

	% -----
	\subsubsection{Section 28:11} % *DC28-11 
	% -----
		\label{DC28-11}

		And again, thou shalt take thy brother, Hiram Page, between him and thee alone, and tell him that those things which he hath written from that stone are not of me and that Satan deceiveth him;

		% \begin{framed}
		% \scriptsize
			% \begin{compactdesc}
				% \item[JSP] 
				% % \item[BC] 
				% % \item[]
			% \end{compactdesc}
		% \end{framed}

	% -----
	\subsubsection{Section 28:12} % *DC28-12 
	% -----
		\label{DC28-12}

		For, behold, these things have not been appointed unto him, neither shall anything be appointed unto any of this church contrary to the church covenants.

		% \begin{framed}
		% \scriptsize
			% \begin{compactdesc}
				% \item[JSP] 
				% % \item[BC] 
				% % \item[]
			% \end{compactdesc}
		% \end{framed}

	% -----
	\subsubsection{Section 28:13} % *DC28-13 
	% -----
		\label{DC28-13}

		For all things must be done in order, and by common consent in the church, by the prayer of faith.

		% \begin{framed}
		% \scriptsize
			% \begin{compactdesc}
				% \item[JSP] 
				% % \item[BC] 
				% % \item[]
			% \end{compactdesc}
		% \end{framed}

	% -----
	\subsubsection{Section 28:14} % *DC28-14 
	% -----
		\label{DC28-14}

		And thou shalt assist to settle all these things, according to the covenants of the church, before thou shalt take thy journey among the Lamanites.

		% \begin{framed}
		% \scriptsize
			% \begin{compactdesc}
				% \item[JSP] 
				% % \item[BC] 
				% % \item[]
			% \end{compactdesc}
		% \end{framed}

	% -----
	\subsubsection{Section 28:15} % *DC28-15 
	% -----
		\label{DC28-15}

		And it shall be given thee from the time thou shalt go, until the time thou shalt return, what thou shalt do.

		% \begin{framed}
		% \scriptsize
			% \begin{compactdesc}
				% \item[JSP] 
				% % \item[BC] 
				% % \item[]
			% \end{compactdesc}
		% \end{framed}

	% -----
	\subsubsection{Section 28:16} % *DC28-16 
	% -----
		\label{DC28-16}

		And thou must open thy mouth at all times, declaring my gospel with the sound of rejoicing.  Amen.

		% \begin{framed}
		% \scriptsize
			% \begin{compactdesc}
				% \item[JSP] 
				% % \item[BC] 
				% % \item[]
			% \end{compactdesc}
		% \end{framed}